% MA 214 Discussion 1 -- covers Lecture 1.
%
% Student handout. Page 1 is the concept review; the question set runs from page 2.
% Compile from inside this folder:  pdflatex Discussion1.tex

\documentclass[11pt]{article}
\usepackage[margin=1in]{geometry}
\usepackage{parskip}
\usepackage{fancyhdr}
\usepackage{array}
\usepackage{booktabs}
\usepackage{enumitem}
\usepackage{amsmath}
\usepackage{tabularx}

\pagestyle{fancy}
\fancyhf{}
\cfoot{\thepage}
\rhead{Discussion 1}
\lhead{MA 214: Applied Statistics}
\renewcommand{\headrulewidth}{.4mm}

\newcommand{\blankline}[1]{\underline{\hspace{#1}}}
\newcolumntype{L}{>{\raggedright\arraybackslash}X}

% Section header: bold title over a hairline rule.
\newcommand{\parttitle}[1]{%
  \par\medskip\noindent
  {\large\bfseries #1}\par
  \vspace{-.5em}\noindent\rule{\linewidth}{0.4pt}\par\smallskip}

% Writing space.
\newcommand{\answerspace}[1]{\vspace{#1}}

% Flush-left question list: the label runs in at the margin and the body wraps
% back to the margin, so nothing is pushed far to the right.
\newlist{questions}{enumerate}{1}
\setlist[questions]{label=\textbf{Question \arabic*.}, wide=0pt, itemsep=1.4em, topsep=.8em}

\newlist{parts}{enumerate}{1}
\setlist[parts]{label=\textbf{(\Alph*)}, leftmargin=1.6em, labelsep=.45em,
  itemsep=.7em, topsep=.5em}

\begin{document}

\noindent\textbf{Name:}\blankline{2.3in}\hfill\textbf{BUID:}\blankline{1.6in}

\begin{center}
  {\Large\bfseries Discussion 1: Randomness, Data, and the Language of Inference}
\end{center}

\noindent\textbf{Goals:} \textbf{(1)} use the notation of population and sample correctly,
\textbf{(2)} say what the sampling distribution is a distribution \emph{of}, \textbf{(3)}
separate what a probability claims from what one sample shows, \textbf{(4)} tell an
observational study from an experiment, \textbf{(5)} recognize the limits of a sampling frame.

\parttitle{Part 1: Concept Review}

{\small

\noindent Everything in Part 2 can be answered from this page. Keep it in front of you.

\begin{enumerate}[label=\textbf{\arabic*.}, wide=0pt, itemsep=.55em, topsep=.5em]

\item \textbf{Random process, outcome, event.} A \textbf{random process} is a repeatable action
whose result is not determined in advance. An \textbf{outcome} is one complete result of running
it once. An \textbf{event} is a \emph{set} of outcomes, which is why a probability attaches to it.
For one bike-share trip: the process is taking a trip, an outcome is one trip with all its
recorded values, and ``the trip lasted over 30 minutes'' is an event.

\item \textbf{Population and sample.} A \textbf{parameter} describes the population and is fixed
but unknown. A \textbf{statistic} is computed from the sample: known, but it changes with every
new sample. You only ever get to see the statistic. An \textbf{inference} is a claim about a
parameter made using a statistic.

\begin{center}
\renewcommand{\arraystretch}{1.1}
\begin{tabular}{lcccc}
\toprule
 & Mean & SD & Proportion & Size \\
\midrule
Population (parameter) & $\mu$ & $\sigma$ & $p$ & $N$ \\
Sample (statistic) & $\bar{x}$ & $s$ & $\hat{p}$ & $n$ \\
\bottomrule
\end{tabular}
\end{center}

\item \textbf{Probability facts.} For an event $A$, and for independent events $A_1,\dots,A_k$,
\[
P(\text{not } A) = 1 - P(A),
\qquad
P(A_1 \text{ and } \cdots \text{ and } A_k) = \prod_{i=1}^{k} P(A_i).
\]
The \textbf{law of large numbers} says that as $n$ grows, $\hat{p} \to p$ and $\bar{x} \to \mu$.
It promises the statistic \emph{settles toward} the parameter, not that it equals the parameter
at any particular $n$.

\item \textbf{Sampling variability and the sampling distribution.} Draw a new sample and every
statistic changes: that is \textbf{sampling variability}. The \textbf{sampling distribution} of
$\bar{x}$ is the distribution of $\bar{x}$ across all possible samples of size $n$ from the
population. One draw from it is one \emph{whole sample}, not one observation, and it does not
depend on which sample you happened to get. As $n$ grows the sampling distribution gets
\emph{narrower}, while the spread of the data, $\sigma$, does not change at all.

\item \textbf{Observational study or experiment.} In an \textbf{experiment} the researcher
\emph{assigns} the explanatory variable, ideally at random. In an \textbf{observational study}
the units arrive with their values already set. A \textbf{confounder} is a variable associated
with both the explanatory variable and the response, which can produce an association with no
causal link behind it. Randomized assignment breaks that link, which is why only an experiment
supports a causal claim.

\item \textbf{Three traps.} \textbf{(a)} Writing $\mu = 14.2$ when $14.2$ came from a sample; it
is $\bar{x} = 14.2$. \textbf{(b)} Calling the histogram of your $n$ observations ``the sampling
distribution''; that is the distribution of the \emph{data}. \textbf{(c)} Believing a larger $n$
fixes a bad sample. It does not: $n$ controls \textbf{variance}, the sampling frame controls
\textbf{bias}, and no amount of the first cures the second.

\end{enumerate}

}

\newpage

\parttitle{Part 2: Question Set}

\noindent\textbf{Scenario.} A city bike-share system recorded every trip taken last year, about
4.2 million of them. You cannot work with all of them, so you take a random sample of
\textbf{250 trips} and record, for each one, the duration in minutes, whether the rider held a
monthly membership, and the station the trip started from. In your sample: $n = 250$ trips, mean
duration $14.2$ min, standard deviation of duration $9.6$ min, and a proportion $0.68$ holding a
membership.

\begin{questions}

%%%%%%%%%%%%%%%%%%%%%%%%%%%%%%%%%%%%%%
\item \textbf{Notation.} Answer in symbols, not prose.

\begin{parts}
\item Write the symbol for each quantity, for all 4.2 million trips and for your 250, and give
the value where you know one.

\begin{center}
\renewcommand{\arraystretch}{1.6}
\begin{tabularx}{\linewidth}{Lccc}
\toprule
\textbf{Quantity} & \textbf{All 4.2M trips} & \textbf{Your 250 trips} & \textbf{Value known} \\
\midrule
Mean duration & & & \\
SD of duration & & & \\
Proportion with membership & & & \\
Number of units & & & \\
\bottomrule
\end{tabularx}
\end{center}

\item Which quantities are \textbf{fixed but unknown}, and which are \textbf{known but would
change} if you drew a new sample of 250?
\answerspace{4em}

\item A classmate writes: ``We found that $\mu = 14.2$ minutes.'' Rewrite the sentence
correctly, and say in one clause what is wrong with the original.
\answerspace{4em}

\item Write, in symbols, the claim ``the true mean trip duration for the whole system is
somewhere near 14 minutes.'' Is that claim a parameter, a statistic, or an inference?
\answerspace{3.5em}
\end{parts}

%%%%%%%%%%%%%%%%%%%%%%%%%%%%%%%%%%%%%%
\item \textbf{Sampling variability.} Three teaching fellows each independently draw their own
random sample of 250 trips from the same year of data.

\begin{center}
\renewcommand{\arraystretch}{1.2}
\begin{tabular}{lccc}
\toprule
 & \textbf{TF 1} & \textbf{TF 2} & \textbf{TF 3} \\
\midrule
Sample mean duration (min) & 14.2 & 13.6 & 15.1 \\
\bottomrule
\end{tabular}
\end{center}

\begin{parts}
\item Why are the three numbers different? Name the phenomenon.
\answerspace{3em}

\item Mark each row \textbf{varies} or \textbf{fixed} across the three TFs.

\begin{center}
\renewcommand{\arraystretch}{1.6}
\begin{tabularx}{\linewidth}{Lc}
\toprule
 & \textbf{Varies / fixed} \\
\midrule
The true mean duration of all 4.2M trips & \\
The sample mean duration & \\
The sample size & \\
The sampling distribution of the sample mean & \\
\bottomrule
\end{tabularx}
\end{center}

\item Complete the sentence precisely: ``The sampling distribution of the sample mean is the
distribution of \blankline{2.7in}\,.''

Then say what is wrong with each of these two answers:
\begin{enumerate}[label=\arabic*., leftmargin=1.8em, itemsep=.15em, topsep=.3em]
\item ``the distribution of the 250 trip durations.''
\item ``the distribution of the true mean.''
\end{enumerate}
\answerspace{5em}

\item The three TFs pool their work and take \textbf{2500} trips instead of 250. For each of the
following answer \emph{increases}, \emph{decreases}, or \emph{stays about the same}: the sample
mean; the SD of trip durations; the spread of the sampling distribution of the sample mean. Give
a reason for the third.
\answerspace{5em}
\end{parts}

%%%%%%%%%%%%%%%%%%%%%%%%%%%%%%%%%%%%%%
\item \textbf{What a probability claims.} A model for this system reports
\[
P(\text{a randomly chosen trip lasts more than 30 minutes}) = 0.08 .
\]

\begin{parts}
\item Restate this as a long-run frequency claim about many trips.
\answerspace{3em}

\item Of your 250 sampled trips, 26 lasted more than 30 minutes. Is the model refuted? Name the
idea you are using.
\answerspace{4.5em}

\item A rider says: ``My last four trips were all under 30 minutes, so my next one is more
likely to go over.'' What is wrong with this reasoning? Name the error.
\answerspace{4em}

\item Under the model, and assuming trips are independent, write an expression for the
probability that \textbf{none} of the next 3 trips lasts more than 30 minutes. Do not evaluate it.
\answerspace{3em}
\end{parts}

%%%%%%%%%%%%%%%%%%%%%%%%%%%%%%%%%%%%%%
\item \textbf{Observational or experiment.} Someone notices in your sample that \textbf{members
average 11.8 minutes per trip and non-members average 19.3 minutes}, and concludes:
\emph{buying a membership makes people take shorter trips.}

\begin{parts}
\item Is this an observational study or an experiment? How do you know?
\answerspace{3em}

\item Name one plausible confounding variable, and say which direction it would push the
comparison.
\answerspace{4.5em}

\item Design an experiment that \emph{would} support the causal claim. Say what you randomize,
to whom, and what you measure.
\answerspace{4.5em}

\item Could the bike-share company actually run your experiment? If not, what stops it?
\answerspace{4em}
\end{parts}

%%%%%%%%%%%%%%%%%%%%%%%%%%%%%%%%%%%%%%
\item \textbf{What this sample can and cannot answer.} Every trip in your sample started from a
station in the \textbf{downtown core}, because that is where the sampling was done. Write
\textbf{yes} if your sample can address the question and \textbf{no} if it cannot, with a short
reason.

\begin{center}
\renewcommand{\arraystretch}{2}
\begin{tabularx}{\linewidth}{Lcp{2in}}
\toprule
\textbf{Question} & \textbf{Yes / no} & \textbf{Reason} \\
\midrule
Mean trip duration for downtown trips? & & \\
Mean trip duration system-wide? & & \\
Do members take shorter trips than non-members downtown? & & \\
Would adding a suburban station increase ridership? & & \\
\bottomrule
\end{tabularx}
\end{center}

\end{questions}

\end{document}
